\subsection{极限与连续性}

设 X、Y 是两个拓扑空间，f 是 X 到 Y 中的映射，B 是点 \(a \in X\) 在 X 中的邻域滤子。我们不把“\(y \in Y\) 是 f 关于滤子 B 的极限”写作 \(y = \lim_{B} f\)，而使用特殊记号

\[
y = \lim _ {t \in I} f (x),
\]

并称 y 为 f 在点 a 的极限，或称当 x 趋于 a 时 \(f(x)\) 趋于 y。类似地，我们不谈 y 是 f 关于 B 的聚点，而谈 y 是 f 在点 a 的聚点。

考察连续性的定义（§ 2，第 1 目，定义 1）与第 3 目的命题 7，即得：

命题 9. 映射 \(f\)（拓扑空间 \(\mathbf{X}\) 到拓扑空间 \(\mathbf{Y}\) 中的）在点 \(a \in \mathbf{X}\) 连续，当且仅当 \(\lim f(x) = f(a)\)。

推论 1. 设 X、Y 是两个拓扑空间，f 是 X 到 Y 中在点 \(a \in X\) 连续的映射；则对 X 上每个收敛于 a 的滤子基 B，滤子基 \(f(\mathfrak{B})\) 都收敛于 \(f(a)\)。反之，若对 X 上每个收敛于 a 的超滤子 U，超滤子基 \(f(\mathfrak{U})\) 都收敛于 \(f(a)\)，则 f 在 a 连续。

第一个断言是命题 9 的直接推论。为证第二个断言，设 f 在 a 不连续；则存在 \(f(a)\) 在 Y 中的一个邻域 W，使得 \(\overline{f}^{1}(W)\) 不属于 a 在 X 中的邻域滤子 B。于是（§ 6，第 4 目，命题 7）存在细于 B 的超滤子 U，它不含 \(\overline{f}^{1}(W)\)，从而包含其余集 \(A = X - \overline{f}^{1}(W)\)（§ 6，第 4 目，命题 5）；由于 \(f(A) \cap W = \emptyset\)，\(f(\mathfrak{U})\) 不收敛于 \(f(a)\)。

推论 2. 设 g 是集合 Z 到拓扑空间 X 中的映射，它关于 Z 上的滤子 F 有极限 a；则若映射 f: X → Y 在 a 连续，合成 f\(\circ\)g 关于 F 以 f(a) 为极限点。

